\documentclass[11pt]{article}
\usepackage[T1]{fontenc}
\usepackage{lmodern}
\usepackage[margin=1in]{geometry}
\usepackage{amsmath,amssymb,amsthm,mathtools,mathrsfs,booktabs,microtype}
\usepackage[colorlinks=true,linkcolor=blue,citecolor=blue,urlcolor=blue]{hyperref}
\newcommand{\R}{\mathbb R}
\newcommand{\Q}{\mathbb Q}
\newcommand{\Z}{\mathbb Z}
\newcommand{\C}{\mathbb C}
\newcommand{\N}{\mathbb N}
\newcommand{\OO}{\mathcal O}
\newcommand{\Emb}{\operatorname{Emb}}
\newtheorem{theorem}{Theorem}[section]
\newtheorem{lemma}[theorem]{Lemma}
\newtheorem{proposition}[theorem]{Proposition}
\theoremstyle{remark}
\newtheorem{remark}[theorem]{Remark}
\title{The irrationality exponent of\texorpdfstring{\\$\log2+\sqrt[3]{2}$}{ log 2 + cube root of 2} is at most six}
\author{Ryan Matthew Casper}
\date{October 13, 2026}

\begin{document}
\maketitle
\begin{abstract}
We prove that $\log2+\sqrt[3]{2}$ has irrationality exponent at most six.
We adapt our previous affine determinant construction to a cubic field
with one compatible embedding. Clearing a single minor gives an
algebraic integer whose norm compares analytic decay at that embedding
with growth at the other two. The norm contains three copies of the
clearing factor and two coarse determinant bounds. We choose parameters that absorb
this full cost, using the logarithmic interpolation and analytic
estimates from our previous proof. A Lean formalization accompanies
the argument.
\end{abstract}

\section{Introduction and main result}
Let $r$ be the positive real root of $r^3=2$ and put $x=\log2+r$.
The number $x$ is transcendental: the Hermite--Lindemann theorem
\cite{lindemann} gives the transcendence of $\log2$, and adding the
algebraic number $r$ preserves it. We study how well $x$ can be
approximated by rational numbers.

For an irrational real number $y$, write $\mu(y)$ for the supremum of
the positive exponents $\nu$ such that
$0<|y-p/q|<q^{-\nu}$ for infinitely many reduced rational numbers
$p/q$ with $q>0$.

\begin{theorem}\label{thm:main}
For every real $\nu>6$ there exists $Q\in\N$, $Q\ge2$, such that
\begin{equation}\label{eq:main}
 q^{-\nu}\le\left|\log2+\sqrt[3]{2}-\frac pq\right|
 \qquad(p\in\Z,\ q\in\N,\ q\ge Q).
\end{equation}
Consequently, $\mu(x)\le6$.
\end{theorem}

The threshold may depend on $\nu$ and is uniform in $p,q$; the
fractions in \eqref{eq:main} need not be reduced. The proof gives
neither a numerical value of $Q$ nor an estimate at $\nu=6$.

Our previous affine theorem \cite[Theorem 3.1]{affine} treats integral
data in a degree-four number field with two compatible embeddings.
Here we use the same determinant construction in degree three, with
one compatible embedding and slope of absolute value at most one.
We retain the integrality assumptions. The logarithmic interpolation
and analytic estimates come from our previous proof, which builds on
OpenAI's determinant method \cite{openai} and our logarithmic
extensions \cite{quadratic}.

The change is in the norm comparison. All three embeddings contribute
to the norm, but only one supplies the analytic decay we need. We bound
the other two coarsely and choose parameters that absorb both bounds.
This works only if we keep the same minor at every embedding: the norm
must belong to one algebraic integer. Section~\ref{sec:estimates}
sets out the three estimates, Lemma~\ref{lem:lift} supplies that common
algebraic integer, and Section~\ref{sec:parameters} chooses the
parameters that make the estimates incompatible.

Rational approximation to $x$ can also be viewed as approximation to
$\log2$ by the cubic numbers $p/q-r$. This explains why a bound for
rational approximation to $\log2$ alone does not give
Theorem~\ref{thm:main}. Section~\ref{sec:context} relates the two height
conventions.

\section{A cubic affine theorem}\label{sec:general}
For a number field $F$, write
$\Emb(F)=\operatorname{Hom}_{\Q\text{-alg}}(F,\C)$ and let
$\OO_F$ be its ring of integers.

\begin{theorem}\label{thm:general}
Suppose $[F:\Q]=3$ and $a,u,v\in\OO_F$. Let $\sigma_0\in\Emb(F)$
have $\sigma_0(a)=b_0\in\R$, where $b_0>0$ and $b_0\ne1$.
Let $y\in\R\setminus\{0\}$, and suppose an embedding
$\tau\in\Emb(F)$ satisfies
\begin{equation}\label{eq:compatible}
 \exp\bigl(\tau(u)(y-\tau(v))\bigr)=\tau(a),
 \qquad |\tau(u)|\le1.
\end{equation}
Then for every real $\nu>6$ there is $Q\in\N$, $Q\ge2$, such that
$q^{-\nu}\le|y-p/q|$ for all integers $p$ and natural $q\ge Q$.
\end{theorem}

The selected embedding $\tau$ need not equal $\sigma_0$. The other two
embeddings are bounded coarsely, whether or not either happens to
satisfy the same exponential identity. The roles are different:
$\sigma_0$ gives a nonzero minor by interpolation, while $\tau$ gives
the analytic upper bound on its determinant.

\subsection{Specialization to the target}
Take $F=\Q(r)$, $a=2$, $u=1$, $v=r$, and $y=x=\log2+r>0$.
Eisenstein's criterion at 2 proves that $T^3-2$ is irreducible, so
$[F:\Q]=3$. The element $r$ is integral, as are $a,u$.
The real embedding is a permissible choice for both $\sigma_0$ and
$\tau$, since $\exp(x-r)=2$.

If $\zeta$ is a primitive cube root of unity, the other
embeddings send $r$ to $r\zeta$ and $r\zeta^2$. They obey
\[
 |\exp(x-r\zeta)|=|\exp(x-r\zeta^2)|
 =2\exp(3r/2)>2.
\]
Thus they do not supply the identity used at the real embedding,
even though every embedding sends the base $a$ to 2.
Theorem~\ref{thm:general} now gives \eqref{eq:main};
Section~\ref{sec:endpoint} deduces the exponent bound.

\section{The inherited matrix estimates}\label{sec:estimates}
We first record the estimates we will use from our previous proof
\cite[Sections 4--8]{affine}. The parameters remain fixed as the matrix
height $H$ tends to infinity. We choose them in
Section~\ref{sec:parameters}, after seeing the full cost of the cubic
norm.

\subsection{Parameters and the selected minor}
Fix $\nu>6$. Assume, for contradiction, that there are approximations
\begin{equation}\label{eq:bad}
 |y-p/q|\le q^{-\nu}
\end{equation}
with arbitrarily large denominators. We select
a finite tuple $p_i/q_i$ with $q_i\ge2$ and $p_i\ne0$, and set
$w_i=\lceil\log q_i\rceil$, $w_* =\min_i w_i$.
The nonzero-numerator requirement causes no loss because $y\ne0$.
Each weight is chosen large enough, after the preceding weights, to
satisfy the separation conditions for logarithmic interpolation
\cite[Sections 4--5]{affine}. Arbitrarily large denominators let us
make these successive choices.

We take $0<\eta<1$ and rational shape parameters satisfying
\begin{equation}\label{eq:shape}
 0<\theta<A<B<1,\quad C>1,\quad CB<1,
 \quad B<C\theta<1,
\end{equation}
For a positive scale $F_0>2/\theta$ and dimension
$m\ge1$, set
\begin{equation}\label{eq:dimension}
 K=\lfloor C^m\rfloor,\qquad w_0=B^{-m},\qquad
 v_0=2K\theta^m w_0.
\end{equation}
The strict volume ratio is $Kw_0\theta^m/v_0=1/2$.

The main quantities in the argument have the following roles.
\begin{center}
\small
\begin{tabular}{@{}lp{0.65\textwidth}@{}}
\toprule
Quantity & Role \\
\midrule
$m,\ H$ & Number of approximations and matrix height \\
$K,\ M(H)$ & Number of multiplicative centers and number of rows \\
$w_i,\ w_0,\ v_0$ & Approximation weights and the weights of $h$ and $s$ \\
$b_H$ & Normalized mean row weight saved by denominator clearing \\
$\mathscr C_H,\ \mathcal C_H$ & Rational clearing factor and analytic collision parameter \\
$E,\ U,\ V$ & Arithmetic error, coarse growth, and analytic error \\
\bottomrule
\end{tabular}
\end{center}

For $H>0$, columns are $(h,\gamma)\in\N\times\N^m$ satisfying
$w_0h+\sum_iw_i\gamma_i\le H$. Rows are $(j,s,\beta)$ with
$0\le j<K$ and $v_0s+\sum_iw_i\beta_i/\theta<H$.
Write $M=M(H)>0$ for the row count. Put
\[
 T_i=\lceil F_0w_i/v_0\rceil,\qquad
 L_i(t)=\sum_{1\le k<T_i}\frac{(-1)^{k+1}}k t^k,
 \qquad c_{\sigma i}=\sigma(u)(p_i/q_i-\sigma(v)).
\]
The matrix at an embedding $\sigma$ is
\begin{equation}\label{eq:matrix}
 \mathcal M_\sigma(H)_{(j,s,\beta),(h,\gamma)}
 =\sigma(a)^{jh}[t^sZ^\beta](1+t)^h
     \prod_{i=1}^m(Z_i+jc_{\sigma i}+L_i(t))^{\gamma_i}.
\end{equation}
At $\sigma_0$, the multiplicative centers $b_0^j$ are distinct and
nonzero. The interpolation theorem allows arbitrary additive centers,
so it gives full row rank at a cofinal set of heights. At each such
height we select one nonzero full-row minor and use exactly those
columns at every embedding. Write $\Delta_\sigma$ for the resulting
determinant.
Define the mean saved row weight by
\[
 b_H=\frac1{MH}\sum_{\text{rows }(j,s,\beta)}\sum_iw_i\beta_i,
 \qquad 0\le b_H\le\theta.
\]

\subsection{Clearing, growth, and decay}
Choose fixed bounds $\kappa,R\ge0$ and $\rho\ge100$ such that
\begin{align*}
 \log\max(1,|\sigma(a)|)&\le\kappa,\\
 |\sigma(u)|(|y|+1+|\sigma(v)|)&\le R
       &&(\sigma\in\Emb(F)),\\
 100|y|&\le\rho,\qquad
 100|\tau(u)(y-\tau(v))|\le\rho.
\end{align*}
There are finitely many embeddings, so we can fix these bounds
before choosing the dimension or approximations. For $x=\log2+r$,
we may take $\kappa=\log2$, $R=6$, $\rho=300$.
Put $\Lambda=\log4+4$ and $c=\log2/4$, as in our previous estimates,
and define
\begin{align}
 E&=\frac{\Lambda F_0m}{v_0}+\Lambda\sum_i\frac1{w_i}
        +\frac\theta{w_*},\label{eq:errors}\\
 U&=\frac{\log2}{v_0}+\frac{\kappa K+\log(3/2)}{w_0}
        +\frac{\log(2(KR+2))}{w_*},\notag\\
 V&=\frac\nu{F_0}+\frac{2\log2}{v_0}+\frac{\rho K}{w_0}
       +\frac{\log4+\log(2K)+\nu+\log(2\rho K)}{w_*}.\notag
\end{align}

The first estimate controls denominator clearing. It supplies a positive rational factor
$\mathscr C_H$ with
\begin{equation}\label{eq:clearcost}
 \frac{\log\mathscr C_H}{MH}\le1-b_H+E.
\end{equation}
The second bounds growth at each of the other two embeddings:
for $\sigma\ne\tau$,
\begin{equation}\label{eq:coarse}
 \frac{\log|\Delta_\sigma|}{MH}\le U+r_c(H),
 \qquad r_c(H)\longrightarrow0.
\end{equation}
At the selected embedding write $\omega=\tau(u)(y-\tau(v))$.
The error comparison is
\[
 |c_{\tau i}-\omega|
 =|\tau(u)|\,|p_i/q_i-y|\le q_i^{-\nu}.
\]
The slope bound therefore lets us use the third estimate without
enlarging the approximation error:
\begin{equation}\label{eq:good}
 \frac{\log|\Delta_\tau|}{MH}
 \le V+r_a(H)+\max\{-\mathcal C_H,-\nu(A(1-\eta)-b_H)\},
 \qquad r_a(H)\longrightarrow0,
\end{equation}
where
\begin{equation}\label{eq:collision}
 \mathcal C_H\longrightarrow\mathcal C_\infty
 =\frac{c\eta^2K\theta^m}{(m+1)v_0A^m}
 =\frac{c\eta^2}{2(m+1)}(B/A)^m.
\end{equation}
These are the estimates of our previous affine construction
\cite[Sections 6--8]{affine}, applied to the same selected minor.
Notice that the analytic frequency $\omega$ and the number $y$ being
approximated need not coincide. In our concrete example they are
$\log2$ and $\log2+r$, respectively. The identity at $\tau$ is enough; we do not take
logarithms at the other embeddings.

\section{The norm of the common minor}\label{sec:norm}
The affine numerators
$n_i=u(p_i-q_iv)$ are integral. Set
$D_i=\operatorname{lcm}(1,\ldots,T_i)$ and
$e_i=\lfloor H/w_i\rfloor$. Multiply each row by
$\prod_iD_i^{e_i}q_i^{-\beta_i}$ and each selected column by
$\prod_iq_i^{\gamma_i}$. Their product is $\mathscr C_H$.

\begin{lemma}\label{lem:lift}
There is one nonzero $Z_H\in\OO_F$ such that
\[
 \sigma(Z_H)=\mathscr C_H\Delta_\sigma
 \qquad(\sigma\in\Emb(F)).
\]
Every selected determinant is consequently nonzero.
\end{lemma}
\begin{proof}
If some $\beta_i>\gamma_i$, the matrix entry is zero. Otherwise put
$\delta_i=\gamma_i-\beta_i$. After coefficient extraction and clearing,
the entry is the image of
\[
 a^{jh}\prod_i\binom{\gamma_i}{\beta_i}
 [t^s](1+t)^h\prod_i D_i^{e_i}(jn_i+q_iL_i(t))^{\delta_i}.
\]
The polynomial $D_i(jn_i+q_iL_i)$ has integral coefficients, and
$\delta_i\le\gamma_i\le e_i$. Hence this entry is integral. These
entries define one matrix over $F$; let $Z_H$ be its determinant.
Determinants commute with every embedding, proving the identity.
Its image at $\sigma_0$ is nonzero, and every embedding is injective.
\end{proof}

We can now use the norm of $Z_H$, which is a nonzero integer. Counting
all three embeddings into $\C$, with each member of any nonreal
conjugate pair counted separately, we obtain
\[
 0\le3\log\mathscr C_H+\sum_{\sigma\in\Emb(F)}\log|\Delta_\sigma|.
\]
Using \eqref{eq:clearcost} and the two bounds \eqref{eq:coarse} gives
\begin{equation}\label{eq:lower}
 \frac{\log|\Delta_\tau|}{MH}
 \ge-3(1-b_H)-3E-2U-2r_c(H).
\end{equation}
This is where degree three enters the proof. The clearing factor
appears three times in the norm, and each of the two other embeddings
contributes one coarse bound. A conjugate pair may have equal absolute
values, but both factors still occur in the product.

\section{Stronger parameters and the contradiction}\label{sec:parameters}
For every $\nu>6$, the shape can be chosen so that
\begin{equation}\label{eq:gap}
 g=\nu\bigl(A(1-\eta)-\theta\bigr)-3(1-\theta)>0.
\end{equation}
For example, take rational $t<1$ close enough to one that
$\nu t>3(1+t)$, put $\theta=t^2$, choose rational
$\theta<A<B<t$ sufficiently close to $t$, then choose
$B/\theta<C<1/B$ and sufficiently small $\eta>0$.
These choices satisfy \eqref{eq:shape} and \eqref{eq:gap}.

Let $S=3E+2U+V$. We need $S<g$ to handle the approximation term in
\eqref{eq:good}, and $3+S<\mathcal C_\infty$ to handle the collision
term. Expanding \eqref{eq:errors} gives
\begin{align}\label{eq:budget}
 S={}&\frac\nu{F_0}
 +\frac{3\Lambda F_0m+4\log2}{v_0}
 +\frac{(\rho+2\kappa)K+2\log(3/2)}{w_0}\notag\\
 &+3\Lambda\sum_i\frac1{w_i}
 +\frac{3\theta+\log4+\log(2K)+\nu+\log(2\rho K)
             +2\log(2(KR+2))}{w_*}.
\end{align}
Choose $0<\epsilon<g$ with $\epsilon\le1/2$, and next choose
$F_0>2/\theta$ with $\nu/F_0<\epsilon/3$.
As $m\to\infty$ in \eqref{eq:dimension},
\[
 m/v_0\to0,\qquad 1/v_0\to0,\qquad K/w_0\to0,
 \qquad 1/w_0\to0,\qquad\mathcal C_\infty\to\infty.
\]
The limits involving $v_0$ use $C\theta/B>1$, the ratio $K/w_0$ tends
to zero because $CB<1$, and the collision parameter grows because
$B/A>1$. Choose one dimension making the two dimension terms of
\eqref{eq:budget} less than $\epsilon/3$ and
$\mathcal C_\infty>4$. Choose a positive rational geometric slack
satisfying the interpolation conditions in \cite[Section 4]{affine}.
The strict volume ratio $1/2$, $K\theta^m<1$, and $\theta<1$ permit
this choice. Then choose a sufficiently large initial
weight threshold so that the remaining inverse-weight terms sum to
less than $\epsilon/3$. Subsequent approximations can still be selected
with the separated-weight inequalities of the same section because
their denominators are unbounded under \eqref{eq:bad}. All these choices
precede the limit in $H$.

We have obtained
\begin{equation}\label{eq:strict}
 S<\epsilon<g,\qquad 3+S<4<\mathcal C_\infty.
\end{equation}
Set $S_H=S+r_a(H)+2r_c(H)$. For all sufficiently large $H$,
\eqref{eq:strict} gives $S_H<g$ and $3+S_H<\mathcal C_H$.
Combining \eqref{eq:good} and \eqref{eq:lower} yields
\[
 -3(1-b_H)\le S_H+
 \max\{-\mathcal C_H,-\nu(A(1-\eta)-b_H)\}.
\]
If the maximum is the approximation term, then
\[
 g+(\nu-3)(\theta-b_H)
 =\nu(A(1-\eta)-b_H)-3(1-b_H)\le S_H<g,
\]
a contradiction. If it is the collision term, then
$\mathcal C_H\le3(1-b_H)+S_H\le3+S_H$, also a contradiction.
Cofinality of full-row-rank heights permits the same selected-minor
construction beyond this common threshold. This proves
Theorem~\ref{thm:general}.

\subsection{Endpoint consequences}\label{sec:endpoint}
The contradiction excludes \eqref{eq:bad} at all sufficiently large
denominators, giving \eqref{eq:main}. The lower bound also supplies
irrationality directly: a rational $x$ would have exact representations
with arbitrarily large unreduced denominators.
For any exponent above six, reduced fractions satisfying its strict
upper bound have bounded denominators, and their numerators are
bounded as well. Thus there are only finitely many such fractions,
which gives $\mu(x)\le6$.

\subsection{Why six appears}
The conditions $B<C\theta$ and $CB<1$ imply $B^2<\theta$, hence
$A<\sqrt\theta$. Consequently
\[
 \frac{A(1-\eta)-\theta}{1-\theta}
 <\frac{\sqrt\theta}{1+\sqrt\theta}<\frac12.
\]
The sufficient gap \eqref{eq:gap} therefore requires $\nu>6$.
Six is the infimum cutoff for this choice of shape using only
$0\le b_H\le\theta$. To lower it, we would need a different shape
or a stronger estimate for the saved row weight. The argument does
not determine the exact value of $\mu(x)$.

\section{Approximation by cubic numbers}\label{sec:context}
Marcovecchio's rational approximation and nonquadraticity estimates
for $\log2$ \cite{marcovecchio} concern different families of
approximants. In our problem the algebraic shift gives a particular
family of cubic numbers.
For a reduced fraction $p/q$, our approximation error is
\[
 |x-p/q|=|\log2-(p/q-r)|.
\]
The cubic number $p/q-r$ has primitive irreducible integer polynomial
\[
 (qT-p)^3+2q^3
 =q^3T^3-3pq^2T^2+3p^2qT-p^3+2q^3.
\]
Irreducibility follows from the rational affine change of variable
from $T^3-2$. For primitivity, any prime dividing $q$ fails to divide
the constant coefficient. Near the fixed target, the coefficient
height is bounded above and below by positive fixed multiples of
$q^3$. Thus a general degree-three approximation measure expressed in
that height would require a factor-three conversion of its exponent
to this rational-denominator convention.

Thus the relevant comparison is with degree-three approximation to
$\log2$ along the family $p/q-r$, after converting its coefficient
height to the denominator $q$. The rational and quadratic estimates
in \cite{marcovecchio} do not give this comparison directly.

\appendix
\section{Formalization}
The accompanying Lean formalization, \texttt{AffineCubic}, proves
Theorems~\ref{thm:main} and~\ref{thm:general}. It constructs the real
field $\Q(2^{1/3})$ and proves the irreducibility, degree, integrality,
and distinguished-embedding identity needed to apply
Theorem~\ref{thm:general}.

The concrete results are in the \texttt{Pilot} module. Alongside the
eventual lower bound and exponent inequality, it proves irrationality
and $\neg\texttt{LiouvilleWith}(\nu,x)$ for every $\nu>6$, using
Mathlib's predicate. The last consequence follows by applying the
lower bound at $(6+\nu)/2$, which absorbs the fixed constant allowed
in that predicate. The companion \path{theorem-map.json} identifies
the declarations, and \path{VERIFY.md} gives the source versions,
reproduction instructions, and proof-checking results.

\section*{Acknowledgements}
OpenAI Codex agents assisted substantially with mathematical development, Lean
formalization, manuscript preparation, and review.

\begin{thebibliography}{9}
\raggedright
\bibitem{affine}
R. M. Casper,
\emph{\href{https://github.com/Mattie/math/blob/cf1338bcee8fe9f08422ec36e8fdf6dc6392c677/preprints/The-irrationality-exponents-of-log-2-plus-sqrt2-and-sqrt2-log-2-are-at-most-4-October-11-2026/paper.pdf}{The irrationality exponents of $\log2+\sqrt2$
and $\sqrt2\log2$ are at most four}}, preprint, October 11, 2026.

\bibitem{openai}
OpenAI,
\emph{\href{https://github.com/openai/math/tree/adc7f1241b42e322a6451854ab7e4b4c146bf78a/preprints/The-irrationality-exponent-of-pi-is-2-September-24-2026}{The irrationality exponent of $\pi$ is $2$}},
preprint, September 24, 2026.

\bibitem{quadratic}
R. M. Casper,
\emph{\href{https://github.com/Mattie/math/tree/4aa817179b5a81358dcdefeb1d454a31ab2f2c27/preprints/The-irrationality-exponent-of-positive-quadratic-logarithms-is-at-most-4-October-8-2026}{The irrationality exponent of positive quadratic
logarithms is at most 4}}, preprint, October 8, 2026.

\bibitem{marcovecchio}
R. Marcovecchio, \emph{The Rhin--Viola method for $\log2$},
Acta Arithmetica \textbf{139} (2009), no.~2, 147--184.
\href{https://doi.org/10.4064/aa139-2-5}{doi:10.4064/aa139-2-5}.

\bibitem{lindemann}
F. Lindemann, \emph{Ueber die Zahl $\pi$}, Mathematische Annalen
\textbf{20} (1882), 213--225.
\href{https://doi.org/10.1007/BF01446522}{doi:10.1007/BF01446522}.
\end{thebibliography}
\end{document}
